% ============================================================
% FILE: latex_additions_v2.tex
%
% Integration instructions:
%
%   BLOCK A — Chapter 3 Architecture Concepts
%     Insert under a new subsection within the architecture concepts chapter:
%       \subsection{Attacks, Hazards and Mitigation Methodologies}
%         \subsubsection{Meltdown and Spectre}
%     Place after existing Section 3.0.8 (Register Concept) and before 3.0.9
%     (Power Concept), or at whichever structural position is appropriate.
%
%   BLOCK B — Chapter 7 (new chapter): Attacks, Hazards and Mitigation
%   Methodologies
%     Create a new chapter file and add the \section{Meltdown and Spectre}
%     block as the first section.
%
%   BLOCK C — Bibliography entries
%     Paste into references.bib.
%
%   ASSUMED ARCHITECTURE EXTENSIONS (for this draft):
%     The content is written assuming the design has been extended beyond the
%     current RV64I single-core baseline to include:
%       - Multi-core (dual- or quad-core) configuration with shared L2 cache
%       - A deeper, explicitly pipelined implementation
%         (Fetch -> Decode -> Issue -> Execute -> Writeback, >= 5 stages)
%       - A shared L2 set-associative data cache in addition to per-core L1
%     These extensions are treated as assumed; when the final specification
%     confirms the exact parameters the placeholder values should be updated.
% ============================================================


% ============================================================
% BLOCK A — Chapter 3 subsection (insert into existing Chapter 3)
% ============================================================

\subsection{Attacks, Hazards and Mitigation Methodologies}
\label{sec:arch_hazards}

\subsubsection{Meltdown and Spectre}
\label{sec:arch_meltdown_spectre}

The extension of the RWU64IM design to a multi-core configuration with a shared
\mbox{L2} cache and a deeper execution pipeline introduces microarchitectural
conditions relevant to transient execution attacks.  These attacks — specifically
Meltdown (CVE-2017-5754)~\cite{lipp2018meltdown} and Spectre
(CVE-2017-5753, CVE-2017-5715)~\cite{kocher2019spectre} — exploit the
discrepancy between what a processor commits architecturally and what it executes
microarchitecturally during speculative or out-of-order processing.

Three architectural preconditions must hold simultaneously for a transient execution
attack to be instantiable~\cite{kocher2019spectre,intel2018analysis}:
(i)~speculative or out-of-order execution,
(ii)~a shared and timing-observable cache hierarchy, and
(iii)~a hardware-enforced isolation boundary (additionally required for Meltdown and
Spectre Variant~2).

The multi-core configuration of this design introduces a shared \mbox{L2} cache
accessible by all cores.  Although the processor operates without a virtual memory
system or hardware privilege separation, the shared cache provides a functional
timing side-channel: access latency to a cached line differs measurably from a cold
miss, and this difference is observable via the \texttt{rdcycle} CSR.  The per-core
hardware branch predictor and the Reorder Buffer (ROB) enabling speculative execution
constitute the remaining preconditions relevant to Spectre Variant~1.

The implications for the current design and the applicable countermeasures are
analysed in detail in Chapter~\ref{chap:attacks}.


% ============================================================
% BLOCK B — Chapter 7 \section{Meltdown and Spectre}
% ============================================================
%
% Chapter preamble (place before this section):
%
% \chapter{Attacks, Hazards and Mitigation Methodologies}
% \label{chap:attacks}
%
% \begin{table}[h]
% \centering
% \caption{Revision History -- Chapter 7}
% \begin{tabular}{llll}
% \hline
% Version & Date & Author & Description \\
% \hline
% 1.0 & January 22, 2026 & Y. Sangale & Initial Draft \\
% \hline
% \end{tabular}
% \end{table}
%
% This chapter analyses microarchitectural security hazards applicable to the
% RWU64IM multi-core RISC-V design, covering transient execution attacks,
% their architectural preconditions, applicability to the implemented
% configuration, and the adopted countermeasures.
% ============================================================

\section{Meltdown and Spectre}
\label{sec:meltdown_spectre}

% ------------------------------------------------------------------
\subsection{Background and Threat Model}
% ------------------------------------------------------------------

Meltdown and Spectre are classes of \emph{transient execution attacks}: they
exploit instructions that are executed microarchitecturally but never committed
to the architectural state~\cite{lipp2018meltdown,kocher2019spectre}.  The
resulting microarchitectural state changes — specifically modifications to the
cache hierarchy — persist after the transient instructions are discarded,
providing a covert channel from which secret data can be recovered by timing
measurement.

The threat model assumes an attacker capable of executing unprivileged code on
the system, controlling the inputs to shared code paths, and measuring execution
time via the \texttt{rdcycle} counter.

% ------------------------------------------------------------------
\subsection{Architectural Preconditions}
\label{subsec:preconditions}
% ------------------------------------------------------------------

Attack feasibility is determined by the simultaneous presence of the following
conditions~\cite{kocher2019spectre,intel2018analysis,arm2018cache}:

\begin{enumerate}
  \item \textbf{Speculative / out-of-order (OoO) execution.}
        A Reorder Buffer (ROB) allows the pipeline to execute instructions ahead
        of their architectural validation.  Without this, no transient execution
        window exists.
  \item \textbf{Hardware branch predictor.}
        Required for Spectre: the predictor enables speculative execution along
        a predicted but architecturally incorrect path.  The BTB (Branch Target
        Buffer) is additionally required for Variant~2.
  \item \textbf{Shared, timing-observable cache.}
        Both Meltdown and Spectre encode transiently accessed data into cache
        state.  The encoding is exfiltrated by measuring the access latency
        difference between a cache hit ($\leq 15$ cycles at \SI{80}{\mega\hertz})
        and a cache miss to DRAM ($\geq 150$ cycles)~\cite{kocher2019spectre}.
        A shared \mbox{L2} cache, accessible across all cores, provides a
        persistent and observable side-channel.
  \item \textbf{Hardware-enforced isolation boundary (attack-specific).}
        Meltdown additionally requires a privilege boundary and a deferred
        page-fault mechanism.  Spectre Variant~2 additionally requires separate
        execution contexts sharing predictor state.
\end{enumerate}

Removing any single required precondition is sufficient to invalidate the
corresponding attack class.

% ------------------------------------------------------------------
\subsection{Meltdown (CVE-2017-5754)}
\label{subsec:meltdown}
% ------------------------------------------------------------------

\subsubsection{Mechanism}

Meltdown exploits the temporal gap between the speculative execution of a
faulting memory load and the enforcement of the access-control decision that
generates the fault~\cite{lipp2018meltdown}.  The pipeline issues the load and
continues executing dependent instructions before the MMU has resolved the
permission check.  During this transient window the loaded value — drawn from a
protected memory region — propagates into a cache-encoding instruction:

\begin{verbatim}
/* Step 1: faulting load from privileged address */
unsigned char secret = *(unsigned char *)kernel_addr;

/* Step 2: encode secret into cache state (transient) */
temp = probe_array[secret * 512];
\end{verbatim}

The subsequent page fault flushes the ROB and restores the architectural
register state, but the cache-line fill caused by \texttt{probe\_array} access
is not reverted.  The attacker recovers the secret by measuring access latency
across all 256 probe-array indices (Flush+Reload
technique~\cite{lipp2018meltdown,gruss2016flushflush}):

\begin{verbatim}
uint64_t t0, t1;
for (int i = 0; i < 256; i++) {
    asm volatile ("rdcycle %0" : "=r"(t0));
    volatile uint8_t v = probe_array[i * 512];
    asm volatile ("rdcycle %0" : "=r"(t1));
    if ((t1 - t0) < HIT_THRESHOLD) recovered_byte = i;
}
\end{verbatim}

A placeholder for the Meltdown attack flow diagram is reserved below:

\begin{figure}[h]
\centering
\fbox{\parbox{0.85\textwidth}{\centering\vspace{2cm}
\textit{[Placeholder: Meltdown attack flow ---
Flush $\to$ Faulting Load $\to$ Speculative Encode $\to$ ROB Flush
$\to$ Reload $\to$ Secret Recovery]}
\vspace{2cm}}}
\caption{Meltdown attack execution flow on an OoO processor with virtual memory.}
\label{fig:meltdown_flow}
\end{figure}

\subsubsection{Applicability to the RWU64IM Design}

The RWU64IM implements physical addressing exclusively.  There are no page tables,
no MMU, no \texttt{satp} register, and no page-fault exceptions.  Every load
instruction targets a physical address and succeeds architecturally; no faulting
access is generated and consequently no transient fault window is created.

In a multi-core extension without virtual memory the situation is unchanged: adding
cores and a shared \mbox{L2} cache does not introduce the page-fault mechanism that
Meltdown requires.  The shared \mbox{L2} does increase the persistence and
cross-core observability of the cache side-channel, but this is a necessary condition
for Meltdown only after the other preconditions are satisfied.

\begin{table}[h]
\centering
\caption{Meltdown Precondition Mapping for RWU64IM (Extended Configuration)}
\label{tab:meltdown_precond}
\begin{tabular}{lll}
\hline
\textbf{Required Precondition} & \textbf{Status in RWU64IM} & \textbf{Impact} \\
\hline
OoO execution / ROB           & Present (per core)        & Precondition met \\
Virtual memory / MMU          & Absent                    & Attack impossible \\
Privilege separation (U/S)    & Absent                    & Attack impossible \\
Page-fault mechanism          & Absent                    & Attack impossible \\
Shared cache (L1/L2)          & Present                   & Channel present \\
\hline
\end{tabular}
\end{table}

\textbf{Result:} Meltdown is structurally not applicable to the RWU64IM in any
single-core or multi-core configuration that retains physical addressing and a single
machine-mode privilege level.

% ------------------------------------------------------------------
\subsection{Spectre (CVE-2017-5753 / CVE-2017-5715)}
\label{subsec:spectre}
% ------------------------------------------------------------------

\subsubsection{Mechanism — Variant 1 (Bounds Check Bypass)}

Spectre Variant~1 exploits the branch predictor's training history to induce
speculative execution of an out-of-bounds memory access~\cite{kocher2019spectre}.
The attacker repeatedly invokes a victim function with in-bounds inputs to train
the predictor, then supplies an out-of-bounds index.  The predictor speculatively
enters the guarded block, executing:

\begin{verbatim}
/* Victim code (shared function) */
if (x < array1_size) {
    uint8_t val = array1[x];          /* out-of-bounds if x is malicious */
    temp = probe_array[val * 512];    /* encodes val into cache */
}
\end{verbatim}

The branch subsequently resolves as not-taken; the ROB discards the speculative
results.  The cache-line fill caused by \texttt{probe\_array[val*512]} persists.
The secret byte is recovered by the same Flush+Reload timing measurement as
Meltdown.

\subsubsection{Mechanism — Variant 2 (Branch Target Injection)}

Variant~2 poisons the Branch Target Buffer to redirect an indirect branch in a
victim execution context to an attacker-chosen code gadget~\cite{kocher2019spectre}.
The gadget accesses sensitive data, encodes it into cache state, and the attacker
recovers it via timing.  This attack requires separate, isolated execution contexts
sharing a physical branch predictor.

A placeholder for the Spectre attack flow diagram is reserved below:

\begin{figure}[h]
\centering
\fbox{\parbox{0.85\textwidth}{\centering\vspace{2cm}
\textit{[Placeholder: Spectre Variant~1 attack flow ---
Train Predictor $\to$ Mispredict $\to$ Speculative OoB Access $\to$
Cache Encode $\to$ Rollback $\to$ Timing Measurement]}
\vspace{2cm}}}
\caption{Spectre Variant~1 execution flow: bounds-check bypass via branch misprediction.}
\label{fig:spectre_v1_flow}
\end{figure}

\subsubsection{Applicability to the RWU64IM Design}

Table~\ref{tab:spectre_precond} maps each Spectre precondition to the extended
RWU64IM configuration.

\begin{table}[h]
\centering
\caption{Spectre Precondition Mapping for RWU64IM (Extended Multi-Core Configuration)}
\label{tab:spectre_precond}
\begin{tabular}{llll}
\hline
\textbf{Precondition} & \textbf{V1} & \textbf{V2} & \textbf{Status in RWU64IM} \\
\hline
Branch predictor (per-core)          & Required & Required & Present \\
OoO execution / ROB (per-core)       & Required & Required & Present \\
Shared cache (L2 across cores)       & Required & Required & Present \\
Isolation boundary (HW-enforced)     & Optional\textsuperscript{a} & Required & Absent \\
Separate execution contexts          & No       & Required & Absent \\
Shared BTB across contexts           & No       & Required & Absent\textsuperscript{b} \\
\hline
\multicolumn{4}{l}{\textsuperscript{a} Required for a meaningful attack; see Variant~1 qualification below.} \\
\multicolumn{4}{l}{\textsuperscript{b} All cores share machine-mode; no separate victim context exists.} \\
\end{tabular}
\end{table}

\paragraph{Variant 1 — single trust domain.}
All software on the RWU64IM executes in machine mode with uniform physical memory
access.  In this configuration there is no security boundary to cross; leakage
across a non-existent boundary has no security consequence.

\paragraph{Variant 1 — multi-component software (conditional risk).}
If the system runs co-resident software components (e.g.\ a cooperative \mbox{RTOS}
with multiple tasks sharing the physical address space), a software-defined trust
boundary emerges.  In this scenario Variant~1 is feasible: the hardware branch
predictor can be trained by one task to induce speculative out-of-bounds reads in
another task's execution context.  The shared \mbox{L2} cache amplifies the
side-channel because a cache-line fill in one core's execution remains observable to
all other cores via timing.  \textbf{This is the primary residual risk of the
multi-core extension.}

\paragraph{Variant 2.}
Variant~2 requires a distinct, isolated victim execution context and shared
predictor state across context boundaries (e.g.\ user/kernel or two OS processes).
No such contexts exist in the RWU64IM privilege model; all code shares machine mode.
Variant~2 is not applicable.

% ------------------------------------------------------------------
\subsection{Cache Timing Side-Channel}
\label{subsec:cache_sidechannel}
% ------------------------------------------------------------------

A cache timing side-channel is present independent of Meltdown/Spectre
applicability~\cite{gruss2016flushflush,arm2018cache}.  The conditions are:

\begin{itemize}
  \item Per-core \mbox{L1} data cache with data-dependent fill behaviour.
  \item Shared \mbox{L2} cache accessible by all cores, whose eviction pattern
        reflects cross-core memory access activity.
  \item \texttt{rdcycle} CSR accessible in machine mode, providing cycle-accurate
        timing resolution sufficient to distinguish a cache hit from a miss.
\end{itemize}

The cross-core observability introduced by the shared \mbox{L2} extends the
attack surface relative to the single-core baseline: a cache-line filled by
speculative execution on core~A is observable from core~B through timing, enabling
cross-core Flush+Reload or Prime+Probe measurements~\cite{gruss2016flushflush}.

% ------------------------------------------------------------------
\subsection{Countermeasures}
\label{subsec:countermeasures}
% ------------------------------------------------------------------

\subsubsection{Software Countermeasures}

\paragraph{Speculation barriers (\texttt{fence}).}
Insert a \texttt{fence} instruction immediately after bounds-check branches at
software trust boundaries.  The fence prevents the pipeline from issuing
instructions beyond that point until all prior instructions have been architecturally
committed, collapsing the speculative execution window:

\begin{verbatim}
if (x < size) {
    asm volatile ("fence");
    val = array[x];
}
\end{verbatim}

\paragraph{Index masking (branch elimination).}
Replace conditional bounds checks with bitwise masking for power-of-two array sizes.
This eliminates the conditional branch and removes the branch predictor dependency
entirely:

\begin{verbatim}
x   = x & (size - 1);   /* no branch: no misprediction possible */
val = array[x];
\end{verbatim}

\paragraph{Constant-time access patterns.}
Security-sensitive routines (cryptographic key operations, inter-task shared
buffers) must access data unconditionally and select results using arithmetic masks
rather than data-dependent branches, eliminating data-dependent cache footprints.

\paragraph{Cycle-counter restriction.}
Setting \texttt{mcounteren.CY~=~0} prevents machine-mode code from delegating
\texttt{rdcycle} access to less-privileged contexts (relevant if U-mode is added in
a future revision).  On the current single-privilege design this reduces the
precision available to tasks within a cooperative multitasking environment if the
scheduler explicitly gates the counter.

\subsubsection{Hardware Countermeasures}

\paragraph{L2 cache partitioning.}
The shared \mbox{L2} cache should be partitioned by software component (e.g.\ by
task or core) using cache-way allocation, if supported by the cache controller
implementation.  This prevents cross-component or cross-core eviction-based
observation.  Each partition must be sized to accommodate the working set of the
assigned component.

\paragraph{Speculative cache-fill suppression.}
The strongest hardware mitigation is to suppress cache-line fills from load
instructions that are subsequently flushed from the ROB without architectural
commitment.  This aligns microarchitectural cache state with the architectural
commit boundary, eliminating the persistent side-channel that all transient
execution attacks depend on.  Implementation requires tagging each speculative
load in the ROB with a \emph{speculative} flag and gating the cache-fill write
until the instruction reaches the commit stage:

\begin{figure}[h]
\centering
\fbox{\parbox{0.85\textwidth}{\centering\vspace{2cm}
\textit{[Placeholder: ROB commit-gated cache-fill suppression diagram ---
Speculative Load $\to$ ROB Entry (speculative flag) $\to$
Commit? $\to$ Yes: allow fill / No: suppress fill]}
\vspace{2cm}}}
\caption{Commit-gated cache-fill suppression: cache updates are held until
ROB commit, preventing transient execution from producing observable side-effects.}
\label{fig:cache_suppression}
\end{figure}

\subsubsection{Applicability Summary}

\begin{table}[h]
\centering
\caption{Countermeasure Applicability — RWU64IM Extended Configuration}
\label{tab:countermeasures}
\begin{tabular}{lllll}
\hline
\textbf{Countermeasure} & \textbf{Targets} & \textbf{Type} &
\textbf{Cost} & \textbf{Effectiveness} \\
\hline
\texttt{fence} barriers          & Spectre V1          & SW & Perf.\ stall  & High \\
Index masking                    & Spectre V1          & SW & Negligible    & High \\
Constant-time access             & Cache side-channel  & SW & Moderate      & High \\
\texttt{mcounteren} restriction  & Cache side-channel  & SW & Negligible    & Medium \\
L2 cache way partitioning        & Cache side-channel  & HW & Area          & High \\
Speculative cache-fill suppress. & All classes         & HW & Area/latency  & Very High \\
\hline
\end{tabular}
\end{table}

% ------------------------------------------------------------------
\subsection{Vulnerability Summary}
% ------------------------------------------------------------------

\begin{table}[h]
\centering
\caption{Transient Execution Vulnerability Assessment — RWU64IM Multi-Core Extension}
\label{tab:vuln_summary_ch7}
\begin{tabular}{lll}
\hline
\textbf{Vulnerability} & \textbf{Status} & \textbf{Determining Factor} \\
\hline
Meltdown (CVE-2017-5754)   & Not applicable        &
  No virtual memory; no page-fault mechanism \\
Spectre V1 (CVE-2017-5753) & Conditionally present &
  HW preconditions met; risk conditional on multi-component SW \\
Spectre V2 (CVE-2017-5715) & Not applicable        &
  No separate execution contexts; no isolated victim domain \\
L1 cache timing channel    & Present (per core)    &
  Shared per-core L1; \texttt{rdcycle} accessible \\
L2 cache timing channel    & Present (cross-core)  &
  Shared L2 observable across all cores \\
\hline
\end{tabular}
\end{table}

The design achieves immunity to Meltdown and Spectre~V2 through the \emph{absence of
required architectural preconditions} rather than through active mitigation.  The
primary residual risk is Spectre~V1 in multi-task software deployments and the
cross-core cache timing channel introduced by the shared \mbox{L2}.  Both are
addressable by the software and hardware countermeasures described in
Section~\ref{subsec:countermeasures}.


% ============================================================
% BLOCK C — Bibliography entries (add to references.bib)
% ============================================================
%
% @inproceedings{lipp2018meltdown,
%   author    = {Lipp, Moritz and Schwarz, Michael and Gruss, Daniel and
%                Prescher, Thomas and Haas, Werner and Horn, Jann and
%                Mangard, Stefan and Kocher, Paul and Genkin, Daniel and
%                Yarom, Yuval and Hamburg, Mike},
%   title     = {{Meltdown: Reading Kernel Memory from User Space}},
%   booktitle = {27th USENIX Security Symposium},
%   year      = {2018},
%   pages     = {973--990},
%   publisher = {USENIX Association}
% }
%
% @inproceedings{kocher2019spectre,
%   author    = {Kocher, Paul and Horn, Jann and Fogh, Anders and Genkin, Daniel and
%                Gruss, Daniel and Haas, Werner and Hamburg, Mike and Lipp, Moritz and
%                Mangard, Stefan and Prescher, Thomas and Schwarz, Michael and
%                Yarom, Yuval},
%   title     = {{Spectre Attacks: Exploiting Speculative Execution}},
%   booktitle = {40th IEEE Symposium on Security and Privacy},
%   year      = {2019},
%   pages     = {1--19},
%   publisher = {IEEE}
% }
%
% @inproceedings{gruss2016flushflush,
%   author    = {Gruss, Daniel and Maurice, Cl{\'{e}}mentine and
%                Fogh, Anders and Mangard, Stefan},
%   title     = {{Flush+Flush: A Fast and Stealthy Cache Attack}},
%   booktitle = {Detection of Intrusions and Malware, and Vulnerability
%                Assessment (DIMVA)},
%   year      = {2016},
%   pages     = {279--299},
%   publisher = {Springer}
% }
%
% @techreport{intel2018analysis,
%   author      = {{Intel Corporation}},
%   title       = {{Analysis of Speculative Execution Side Channels}},
%   institution = {Intel Corporation},
%   year        = {2018},
%   note        = {White Paper, Revision~4.0}
% }
%
% @techreport{arm2018cache,
%   author      = {{ARM Limited}},
%   title       = {{Cache Speculation Side-channels}},
%   institution = {ARM Limited},
%   year        = {2018},
%   note        = {Whitepaper, Version~2.4}
% }
%
% @techreport{riscv2022zicbom,
%   author      = {{RISC-V International}},
%   title       = {{RISC-V Cache Management Operation ISA Extensions (Zicbom)}},
%   institution = {RISC-V International},
%   year        = {2022},
%   note        = {Version~1.0}
% }
% ============================================================
